Implementasi \emph{Convolutional Neural Network} (CNN) pada sistem tertanam membutuhkan rancangan yang akurat, rendah latensi, hemat daya, dan dapat terintegrasi dengan jalur akuisisi citra. Penelitian ini merancang dan mengintegrasikan akselerator CNN berbasis komputasi \emph{fixed-point} 16-bit pada board FPGA Digilent Arty A7-100T, dengan kamera OV7670 sebagai sumber citra, DDR3 SDRAM sebagai \emph{frame buffer}, modul \emph{region of interest}, CNN accelerator berbasis Verilog, modul \emph{argmax}, dan keluaran VGA. Model CNN dilatih menggunakan Python pada dataset MNIST berukuran $28 \times 28$ piksel, kemudian bobot dan bias dikuantisasi ke format Q2.14 yang tidak menunjukkan saturasi. Hasil implementasi menunjukkan CNN accelerator menghasilkan keluaran inferensi dalam 3802 siklus clock atau 38,02 $\mu$s pada 100 MHz, dengan throughput sekitar 26302 inferensi per detik. Estimasi daya total on-chip sistem terintegrasi adalah 1,099 W, dan penggunaan DSP mencapai 235 dari 240 DSP yang tersedia atau sekitar 97\% pada board Arty A7-100T, sehingga integrasi kamera, memori, akselerator CNN, dan tampilan dapat direalisasikan dengan batas utama pada margin DSP yang sangat kecil.

\bigskip
\noindent
\textbf{Kata kunci :} \emph{CNN}, \emph{FPGA}, \emph{fixed-point} 16-bit, Arty A7-100T, OV7670, akselerator CNN.